\section{Introduction\label{sec:intro}}

Maintaining a software architecture requires knowing which implementation elements belong to its modules. That correspondence can become difficult to maintain as a system grows: new files introduce responsibilities and dependencies, while architectural documentation may lag behind the implementation. Code-to-architecture (C2A) mapping addresses this problem by assigning source files to the modules of a reference architecture. The resulting mapping supports reflexion modeling, which compares implementation dependencies with architectural relations to identify divergence~\cite{reflexion1}. It therefore connects the modular decomposition used to reason about a system~\cite{parnas1972,clements2003} to the code in which that decomposition is realized.
 
Several sources of evidence can support a mapping. Identifiers express domain concepts and implementation roles. Contextual code representations capture information from declarations and operations. Static dependencies locate a file among other files with which it interacts, while directories record how developers organize the project~\cite{sarif,negar,garcia2013}. These sources need not agree: a utility may depend on several modules, terminology may be shared across responsibilities, and a directory may contain files assigned to different architectural modules.
 
The usefulness of an information source also depends on how it enters a predictor. A dependency embedding can encode graph structure before a classifier sees a file. Propagating over dependencies then supplies a second form of access to related files, with its own assumptions about direction, relation types, and aggregation. A gain from the embedding does not establish that this additional propagation will help. Similarly, a folder-aware model can improve overall agreement with a reference architecture while losing correct predictions for files that depart from a directory-based assignment. These distinctions motivate studying both the information supplied to a model and the changes caused by the way it processes that information.

Prior architecture recovery and mapping work provides several ways to represent these sources. NEGAR learns structural representations using Node2Vec random walks on the dependency graph~\cite{negar}; GAER used a graph autoencoder with Word2Vec identifier embeddings for recovery~\cite{gaer}; prior mapping work uses lexical and organizational features with graph convolutional networks~\cite{jabeen2025graph}. Later work compares random walk-based methods with graph attention and diffusion-based rewiring and examines alternative uses of dependencies and folders~\cite{structsar}. These studies generally fix the representation chosen for each source, so it remains unclear which signal provides the strongest architectural evidence and how it should be represented.

We study this question through C2A mapping, a prerequisite for reflexion modeling~\cite{reflexion1}. We formulate C2A mapping as a classification problem in which each file is assigned to one module of the reference architecture. We work in a low-label setting, where about 5\% of the files are already mapped, and the model infers the remaining assignments. The model is trained either on the labeled files alone or with self-training using its own confident predictions. This setting allows us to examine how file representations and graph relations contribute under limited supervision.

Our evaluation uses ten Java, C, and C++ systems and covers lexical Word2Vec, structural Node2Vec, and contextual UniXcoder representations. These serve as inputs to a graph-free neural baseline and to two graph neural networks (GNNs), a Heterogeneous Graph Transformer (HGT)~\cite{hgt} and, in the controlled comparisons, a Graph Attention Network (GAT)~\cite{gat}. The heterogeneous configurations contain file nodes and, where applicable, folder nodes and one aggregate member node per file. They vary both the context available to a file and where feature content is stored. We compare these file representations and complement the configuration comparisons with controlled comparisons that change one design choice while holding the features, labeled splits, initialization, and training budgets fixed within each comparison.
 
The configuration results motivate two complementary analyses. A directory analysis distinguishes errors corrected by a model from previously correct predictions that it makes incorrect. A dependency analysis first tests whether embeddings retain useful information beyond a degree-preserving rewiring control on three systems, then compares ways of using dependency neighbors on all nine systems. Both stages of the dependency analysis use supervised-only and self-training schedules of equal duration.

We address three research questions:
\begin{itemize}
    \item \textbf{RQ1.} How does mapping performance vary across lexical, structural and contextual file representations?
    \item \textbf{RQ2.} How does mapping performance change when those representations are used in file, member, and folder graph configurations?
    \item \textbf{RQ3.} How do graph representation, training schedule, and neighbor processing affect the contribution of structural context?
\end{itemize}

The results show that dependency information supports C2A mapping through both file representations and neighbor context, but the benefit of neighbor context depends on how it is processed. Folder organization yields the largest gains, largely reflecting information already available through a directory rule. Aggregate improvements can conceal regressions on individual files. These findings clarify how input signals and their use by a predictor affect C2A mapping across the nine systems studied\footnote{\label{fn:code}Code and data are available at \url{https://github.com/sse-lnu/hetmap}.}.